% C_n 的兩個問題 × 兩個機關（1387–1392）。三格灰階，第二問對法院一格用強調色。
\begin{tikzpicture}[qa,cell/.style={nbox,text width=52mm,minimum height=15mm}]
  \node[hdr] at (3.15,1.2) {立法機關（約五百人）};
  \node[hdr] at (8.85,1.2) {聯邦最高法院（九人）};
  \node[hdr,anchor=east,align=right] at (0.15,0) {第一問\\[-1pt]{\scriptsize\mdseries 為何由這些人}};
  \node[hdr,anchor=east,align=right] at (0.15,-1.8) {第二問\\[-1pt]{\scriptsize\mdseries 為何採此程序}};
  \node[cell] at (3.15,0) {公平選舉：C\textsubscript{n} 與全體公民受到平等對待（1387）};
  \node[cell] at (8.85,0) {總統提名、參議院同意的間接民主資格，與議員相較遠為不足（1391）};
  \node[cell] at (3.15,-1.8) {多數決的公平性論證：中立、平等，給每一意見最大可能的權重（1388）};
  \node[cell,kbox] at (8.85,-1.8) {五票勝過四票；大法官不代表任何人，多數決在此欠缺道德基礎（1391--1392）};
\end{tikzpicture}
